\section{Objectives and Patient-Facing Endpoints}

\subsection{What the study is trying to find out}

An objective in a protocol is written as a noun phrase, and a noun phrase is
hard to argue with because it does not commit to anything. Restated as
questions, the objectives of this study are these.

\textbf{Primary.} Is the combination of perioperative daraxonrasib and an
on-premises LLM-directed eight-arm robotic Whipple safe enough, and practical
enough, to be carried forward? Safety is measured three ways: the rate of
dose-limiting toxicity across the three dose levels, the rate of device-related
and procedure-related serious adverse events through day 30, and the rate of
Clavien-Dindo grade III or higher complications through day 30.

\textbf{Secondary.} How did the operation and the drug perform on the outcomes
that matter to the disease? Four measures: the R0 resection rate at day-90
pathology, the ISGPS grade B/C postoperative pancreatic fistula rate
\cite{Bassi2017}, 90-day mortality, and progression-free and overall survival to
24 months.

\textbf{Exploratory.} What did the advisory layer actually contribute? Three
measures: how often the surgeon accepted the model's proposed step, how often
they declined it, and how often the deterministic gate rejected a proposal
before the surgeon saw it. These are reported whether or not they are
favourable, because a system whose override rate is never published cannot be
audited for automation bias \cite{FDATrans2024}.

\subsection{Every endpoint, and the sentence it will let someone say}

The parent protocol stops at the endpoint definition \cite{onpremwhippl}. A
participant who reads only that definition learns that the study measures an
ISGPS grade B/C postoperative pancreatic fistula rate, and is no better informed
than before. Figure 10 carries each endpoint one rank further, to the plain
sentence that endpoint will authorise a clinician to say to the participant or
their family once it reads out.

Three of those sentences are marked apart in the figure, because they are the
three that patients in the surveyed literature ranked highest: the margins came
back clear, you were alive at three months, and the cancer had not returned by
this visit \cite{WuSemiAI2026, TelesAI2026A}. The others matter, but those three
are what people actually asked to be told.

\begin{pafig}[mermaid-type: flowchart TD, three ranks]
\node[mmgoal,text width=34mm] (o1) at (-6.4,0)  {\textbf{Primary objective}\\characterise safety and feasibility of the combination};
\node[mmgoal,text width=34mm] (o2) at (1.4,0)   {\textbf{Secondary objective}\\estimate oncologic and survival performance};
\node[mmstep,text width=34mm] (o3) at (9.2,0)   {\textbf{Exploratory objective}\\quantify what the advisory layer contributed};
\node[mmin,text width=25mm] (p1) at (-9.0,-3.4) {DLT rate across DL1 to DL3};
\node[mmin,text width=25mm] (p2) at (-6.4,-3.4) {Device and procedure SAE rate to day 30};
\node[mmin,text width=25mm] (p3) at (-3.8,-3.4) {Clavien-Dindo grade III or higher};
\node[mmin,text width=25mm] (p4) at (-1.2,-3.4) {R0 resection rate at day-90 pathology};
\node[mmin,text width=25mm] (p5) at (1.4,-3.4)  {ISGPS grade B/C fistula rate};
\node[mmin,text width=25mm] (p6) at (4.0,-3.4)  {90-day mortality};
\node[mmin,text width=25mm] (p7) at (6.6,-3.4)  {PFS and OS to 24 months};
\node[mmgray,text width=25mm](p8) at (9.2,-3.4) {Advisory concordance and override rate};
\node[mmsoft,text width=25mm](m1) at (-9.0,-6.8){``The dose you were given was tolerated at your level.''};
\node[mmsoft,text width=25mm](m2) at (-6.4,-6.8){``Nothing the device did caused you serious harm.''};
\node[mmsoft,text width=25mm](m3) at (-3.8,-6.8){``Your recovery needed no unplanned major intervention.''};
\node[mmmid,text width=25mm] (m4) at (-1.2,-6.8){``The margins came back clear.''};
\node[mmsoft,text width=25mm](m5) at (1.4,-6.8) {``Your pancreas join healed without a draining leak.''};
\node[mmmid,text width=25mm] (m6) at (4.0,-6.8) {``You were alive at three months.''};
\node[mmmid,text width=25mm] (m7) at (6.6,-6.8) {``The cancer had not returned by this visit.''};
\node[mmsoft,text width=25mm](m8) at (9.2,-6.8) {``Your surgeon disagreed with the model here, and was free to.''};
\draw[mmedge] (o1) to[out=-90,in=90,looseness=0.7] (p1);
\draw[mmedge] (o1) -- (p2);
\draw[mmedge] (o1) to[out=-90,in=90,looseness=0.7] (p3);
\draw[mmedge] (o2) to[out=-90,in=90,looseness=0.7] (p4);
\draw[mmedge] (o2) -- (p5);
\draw[mmedge] (o2) to[out=-90,in=90,looseness=0.7] (p6);
\draw[mmedge] (o2) to[out=-90,in=90,looseness=0.7] (p7);
\draw[mmedge] (o3) -- (p8);
\foreach \i in {1,...,8}{\draw[mmedgeb] (p\i) -- (m\i);}
\node[font=\scriptsize\sffamily\itshape,text=protogray,anchor=north west,align=left]
  at (-10.6,-9.6) {Rank three exists in no clinical protocol. It is this paper's contribution: an endpoint that cannot be restated as a sentence about\\the participant has not yet been explained to them. The three lighter-blue sentences are the three the surveyed patients ranked highest.};
\end{pafig}
\vspace{-0.7cm}
\figcaption{Figure 10. Each objective carried through its endpoints to the plain sentence that\\
endpoint will authorise, with the three sentences patients in the surveyed literature\\
ranked highest, clear margins, being alive, and no return, marked apart from the rest.}

\subsection{The endpoint table, with a column the protocol does not carry}

\vspace{0.30em}

\noindent\begin{tabularx}{\textwidth}{L{2.4cm} L{4.1cm} C{2.0cm} Y}
\toprule
\textbf{Class} & \textbf{Endpoint} & \textbf{Timepoint} & \textbf{What it means for you} \\
\midrule
Primary & Dose-limiting toxicity rate, DL1 to DL3 & Day 30 & The dose you were given was tolerated at your level \\
Primary & Device- and procedure-related serious adverse event rate & Day 30 & Nothing the device did caused you serious harm \\
Primary & Clavien-Dindo grade III or higher rate & Day 30 & Your recovery needed no unplanned major intervention \\
Secondary & R0 resection rate & Day 90 & The margins came back clear \\
Secondary & ISGPS grade B/C fistula rate & Day 30 & Your pancreas join healed without a draining leak \\
Secondary & 90-day mortality & Day 90 & You were alive at three months \\
Secondary & Progression-free and overall survival & 24 months & The cancer had not returned by this visit \\
Exploratory & Advisory concordance and override rate & Per procedure & Your surgeon disagreed with the model here, and was free to \\
\bottomrule
\end{tabularx}

\vspace{0.30em}

\subsection{The endpoint registry as a typed record}

An endpoint is not a sentence, it is a record. It has a definition, a timepoint,
an analysis population it is computed over, a class of evidence, and, in this
paper, a plain meaning. Drawing the registry as a schema rather than a list
forces every one of those fields to be filled for every endpoint, and exposes
any endpoint where one is empty. Figure 11 does that, with foreign keys to the
analysis populations and to the four provenance classes that \S 10 uses.

The `plain meaning' row is the one field no clinical protocol carries. Treating
it as a required field of the record, rather than as a nice addition, is the
editorial position of this paper expressed as a schema constraint.

\begin{pafig}[d2-type: sql\_table]
\node[d2sql,text width=48mm] (t1) at (0,0) {%
  \textbf{primary\_safety}\nodepart{two}%
  \textit{endpoint\_id} \hfill varchar PK\\
  dlt\_rate \hfill proportion, DL1 to DL3\\
  device\_sae\_rate \hfill proportion to day 30\\
  clavien\_dindo\_ge3 \hfill proportion to day 30\\
  timepoint \hfill day 30\\
  population \hfill DLT-evaluable FK\\
  \colorbox{pablue2}{plain\_meaning} \hfill the dose was tolerated\\
  provenance\_class \hfill P, protocol-specified};
\node[d2sql,text width=48mm] (t2) at (0,-6.6) {%
  \textbf{secondary\_oncology}\nodepart{two}%
  \textit{endpoint\_id} \hfill varchar PK\\
  r0\_rate \hfill proportion at day 90\\
  isgps\_bc\_fistula \hfill ISGPS 2016 grade\\
  mortality\_90d \hfill proportion at day 90\\
  pfs, os\_24mo \hfill time to event, months\\
  population \hfill safety analysis set FK\\
  \colorbox{pablue2}{plain\_meaning} \hfill margins clear, join held\\
  provenance\_class \hfill M measured, C comparator};
\node[d2sql,text width=48mm] (t3) at (0,-13.0) {%
  \textbf{exploratory\_advisory}\nodepart{two}%
  \textit{endpoint\_id} \hfill varchar PK\\
  advisory\_concordance \hfill proportion accepted\\
  override\_rate \hfill proportion declined\\
  gate\_rejection\_rate \hfill blocked pre-approval\\
  restart\_window\_hit \hfill T+7 / T+14 / T+21\\
  population \hfill per-procedure FK\\
  \colorbox{pablue2}{plain\_meaning} \hfill the surgeon disagreed\\
  provenance\_class \hfill S, simulation anchor};
\node[d2sql,text width=48mm] (t4) at (9.6,-2.6) {%
  \textbf{analysis\_population}\nodepart{two}%
  \textit{population\_id} \hfill varchar PK\\
  safety\_analysis\_set \hfill any intervention received\\
  dlt\_evaluable \hfill completed the DLT window\\
  per\_protocol \hfill no major deviation\\
  itt\_equivalent \hfill all enrolled\\
  \colorbox{pablue2}{who\_you\_are\_in} \hfill the safety set, from your first study procedure onward};
\node[d2sql,text width=48mm] (t5) at (9.6,-9.6) {%
  \textbf{provenance\_class}\nodepart{two}%
  \textit{class\_id} \hfill char PK\\
  M \hfill measured in a cited human series\\
  C \hfill contemporaneous nationwide comparator\\
  S \hfill author simulation, not human data\\
  P \hfill protocol-specified limit, not a result\\
  \colorbox{pablue2}{rule} \hfill every number in \S 10 carries one of these four letters};
\draw[d2edgeb] (t1.east) to[out=0,in=180,looseness=0.85] node[font=\tiny\sffamily,above,pos=0.45]{one to many} (t4.west);
\draw[d2edgeb] (t2.east) to[out=0,in=180,looseness=0.85] (t4.west);
\draw[d2edgeb] (t3.east) to[out=0,in=180,looseness=0.85] (t4.west);
\draw[d2edge] (t1.east) to[out=-25,in=170,looseness=0.9] node[font=\tiny\sffamily,above,pos=0.6]{typed} (t5.west);
\draw[d2edge] (t2.east) to[out=-15,in=175,looseness=0.9] (t5.west);
\draw[d2edge] (t3.east) to[out=15,in=185,looseness=0.9] (t5.west);
\node[d2gray,text width=88mm,anchor=north west,align=left] at (-3.2,-17.4)
  {The \texttt{plain\_meaning} row exists in no clinical protocol. It is the contribution of this paper: an endpoint that cannot be restated as a sentence about the participant has not yet been explained to them.};
\end{pafig}
\vspace{-0.7cm}
\figcaption{Figure 11. The endpoint registry drawn as typed records with foreign keys to the\\
analysis populations and to the provenance classes, and with a plain-meaning row that no\\
clinical protocol carries but that this paper treats as a required field of the record.}

\subsection{What a positive result would and would not establish}

Suppose the study completes all three dose levels with no more than one
dose-limiting toxicity in six at any level, no device-related serious adverse
event through day 30, an R0 rate consistent with the nationwide robotic
comparator, and an ISGPS grade B/C fistula rate at or below the phase-4 figure
from that comparator \cite{DutchCohort2025}. That is the best realistic outcome.

It would establish that the combination can proceed to a study that could test
survival. It would not establish that the combination improves survival, and it
would not establish that the robotic approach is superior to an open operation
for cancer control. Eighteen participants cannot support either claim, the FDA's
own information page states that long-term cancer outcomes have not been
established for every use of robotically assisted surgical devices
\cite{FDARobot2022}, and the National Cancer Institute's guidance says the same
thing about early-phase studies in general \cite{NCITrial2024}.

The reason to say this plainly, in a paper whose purpose is to argue for the
trial, is that the surveyed literature shows patients discount claims that
outrun their evidence and trust the sources that mark their own limits
\cite{Jauniaux2025}. An advocacy document that overstated the primary endpoint
would undermine the twenty-one specific answers the rest of the paper gives.
